\documentclass[11pt]{article}
\usepackage[T1]{fontenc}
\usepackage{lmodern,microtype}
\usepackage[margin=1in]{geometry}
\usepackage{amsmath,amssymb,amsthm,mathtools}
\usepackage{booktabs,array,caption,enumitem,needspace}
\usepackage[hidelinks]{hyperref}
\newtheorem{theorem}{Theorem}[section]
\newtheorem{proposition}[theorem]{Proposition}
\newtheorem{lemma}[theorem]{Lemma}
\theoremstyle{definition}
\newtheorem{example}[theorem]{Example}
\theoremstyle{remark}
\newtheorem{remark}[theorem]{Remark}
\newcommand{\C}{\mathbb C}
\newcommand{\R}{\mathbb R}
\newcommand{\Z}{\mathbb Z}
\newcommand{\tr}{\tau}
\newcommand{\spec}{\operatorname{spec}}
\newcommand{\ord}{\operatorname{ord}}
\newcommand{\PP}{\operatorname{PP}}
\newcommand{\sot}{\mathrm{sot}}
\captionsetup{font=small,labelfont=bf,skip=6pt,hypcap=false}
\setlist[enumerate]{itemsep=3pt}
\setlength{\emergencystretch}{1.5em}
\renewcommand{\arraystretch}{1.16}
\title{Intersections of subfactors in a one-parameter\\Hadamard family}
\author{}
\date{}
\begin{document}
\maketitle
\vspace{-2em}

\section{Introduction}
Two subfactors can each have index \(4\) in the same factor while their
intersection has a much larger index, or even infinite index. This note
examines that phenomenon for the one-parameter family of \(4\times4\)
Hadamard matrices. The guiding questions are: when do two parameters give
the same subfactor, when do their conditional expectations commute, and
what determines the index of their intersection?

The answers supplied by Bakshi--Guin are compared with finite-dimensional
calculations. We first introduce the objects and explain the linear algebra
behind the computations. We then state the relevant results in the
parameters \(s,t\) used here and discuss every pair that was calculated.
The additional \(S(1,1,2,2)\) calculation is treated separately in
Appendix~\ref{sec:ah}, since its parameters are not governed by the
Hadamard classification.

The numerical tables illustrate theorems; they are not proofs of
statements about infinitely many levels. In particular, a sequence which
grows through five levels need not continue growing indefinitely.

\section{The family and the notation}\label{sec:setup}
\subsection{The ambient factor and the intersection}
A \emph{von Neumann algebra} is a unital algebra of bounded operators on a
Hilbert space, closed under adjoints and under strong operator limits.
Strong convergence means convergence on each vector. Its \emph{center}
consists of the elements commuting with every element of the algebra.
A von Neumann algebra whose center is \(\C1\), the scalar multiples of
the identity, is called a \emph{factor}. The symbol \(x^*\) denotes
the adjoint, or the conjugate transpose when \(x\) is a matrix.

Here \(M\) is the hyperfinite type \(\mathrm{II}_1\) factor, with normalized
trace \(\tr\). Thus \(\tr(1)=1\) and \(\tr(xy)=\tr(yx)\); the trace is
faithful, so \(\tr(x^*x)=0\) implies \(x=0\). The term \emph{hyperfinite}
means that \(M\) can be approximated by an increasing sequence of
finite-dimensional algebras. Unlike a matrix algebra, a type
\(\mathrm{II}_1\) factor has no nonzero minimal projections.
A projection is an element \(p=p^*=p^2\); it is minimal if there
is no projection other than \(0,p\) lying below it.
A \emph{subfactor} is a unital von Neumann subalgebra that is itself a factor.

For \(r\in\R\), measured in radians, put
\begin{equation}\label{eq:family}
 H(r)=
 \begin{pmatrix}
 1&1&1&1\\
 1&i e^{ir}&-1&-i e^{ir}\\
 1&-1&1&-1\\
 1&-i e^{ir}&-1&i e^{ir}
 \end{pmatrix},
 \qquad u(r)=\frac12 H(r).
\end{equation}
The entries of \(H(r)\) have modulus one and \(H(r)H(r)^*=4I_4\);
this is the Hadamard condition. Consequently \(u(r)\) is unitary.
The spin-model construction associates to this matrix a subfactor
\(N_r\subset M\) of index \(4\). We use the simultaneous realization
provided by the construction: all the \(N_r\) sit in the \emph{same}
\(M\). This matters, since an intersection depends on the embeddings,
not merely on the abstract isomorphism classes of its two members.

Fix \(s,t\in\R\). Our main object and our relative-phase notation are
\begin{equation}\label{eq:phase}
 P_{s,t}=N_s\cap N_t,\qquad
 \Delta=t-s,\qquad \theta=e^{i\Delta},\qquad
 q=\ord(\theta^2).
\end{equation}
The order of a complex number on the unit circle is the least positive
integer \(j\) for which its \(j\)-th power is \(1\), and is \(\infty\)
if no such integer exists. Thus
\begin{equation}\label{eq:denominator}
 \frac{\Delta}{\pi}=\frac pq,\quad
 \gcd(|p|,q)=1,\ q>0
 \quad\Longrightarrow\quad
 \ord(e^{2i\Delta})=q.
\end{equation}
If \(\Delta/\pi\) is irrational, this order is infinite.
For example, \(\Delta=\pi/3\) gives \(q=3\), whereas \(\Delta=1\)
gives \(q=\infty\). The latter parameter means one \emph{radian}.

\subsection{Index and conditional expectation}
For a von Neumann subalgebra \(P\subset M\), let \(E_P:M\to P\) be
the trace-preserving conditional expectation. It keeps the part of
an operator belonging to \(P\), characterized by
\(\tr(y^*E_P(x))=\tr(y^*x)\) for every \(y\in P\).
On the Hilbert space
\(L^2(M,\tr)\), obtained by completing \(M\) in the norm
\(\|x\|_2=\tr(x^*x)^{1/2}\), this map extends to the orthogonal
projection onto \(L^2(P,\tr)\).

The Jones index \([M:N]\) measures the size of a factor inclusion
\(N\subset M\); formally it is the dimension of \(L^2(M)\) as an
\(N\)-module, with dimension measured using the trace. It is not the
ordinary complex vector-space dimension. To discuss intersections whose
factoriality has not been established, we use the Pimsner--Popa index:
\begin{equation}\label{eq:pp}
 [M:P]_{\PP}
 =\left(\sup\{c\geq0:E_P(x)\geq c x
                    \text{ for every }x\in M,\ x\geq0\}\right)^{-1}.
\end{equation}
Here \(x\geq0\) means that \(x\) is a positive operator, and \(1/0=\infty\).
For inclusions of type \(\mathrm{II}_1\) factors this agrees with the
Jones index. Write \(I(s,t)=[M:P_{s,t}]_{\PP}\).
This notation does not presume that every \(P_{s,t}\) is a factor.

\section{What a finite-level calculation measures}\label{sec:finite}
\subsection{The finite algebras}
The construction gives increasing finite-dimensional algebras
\[
 D_n=A_{n,1}\subset M,\qquad
 B_n(r)=A_{n,0}^{(r)}\subset D_n\quad(n\geq1).
\]
The \(D_n\) do not depend on \(r\); the subalgebras \(B_n(r)\) do.
The notation \(A_{n,j}\) agrees with the implementation; in this note
\(D_n\) and \(B_n(r)\) make the role of each algebra explicit.
Their strong closures recover the infinite algebras:
\[
 M=\overline{\bigcup_nD_n}^{\,\sot},\qquad
 N_r=\overline{\bigcup_nB_n(r)}^{\,\sot}.
\]
For a fixed pair define
\begin{equation}\label{eq:finiteintersection}
 C_n=C_n(s,t)=B_n(s)\cap B_n(t),\qquad
 b_n=\dim_{\C}B_n(s),\qquad c_n=\dim_{\C}C_n.
\end{equation}
In the Hadamard calculation \(b_n=4^n\), for either parameter.
At the first level \(D_1\cong M_4(\C)\) and \(B_1(r)\cong\C^4\).
Hence \(b_1=4\), although the containing algebra has dimension \(16\).
Writing \(\mathcal D_4\) for the diagonal matrices, the starting square
in the convention of this calculation is
\[
 \begin{matrix}
 u(r)^*\mathcal D_4u(r)&\subset&M_4(\C)\\
 \cup&&\cup\\
 \C1&\subset&\mathcal D_4.
 \end{matrix}
\]
The Hadamard condition makes the two expectations onto the
four-dimensional algebras compose to \(x\mapsto\tr(x)1\).
The spin-model construction iterates this square by the Jones basic
construction, adjoining the orthogonal projections associated with the
inclusions. We take that construction as given; the calculations below
use its finite algebras and their trace-preserving inclusions.

The compatibility of conditional expectations in this construction gives
\(E_{N_r}(D_n)=B_n(r)\), and
\begin{equation}\label{eq:closure}
 P_{s,t}=\overline{\bigcup_n C_n(s,t)}^{\,\sot}.
\end{equation}
This is a property of these towers, not a general rule that an
intersection commutes with taking closures. One way to see the reason
is to restrict the two orthogonal projections \(E_{N_s},E_{N_t}\) to
\(L^2(D_n)\). Alternating the projections converges to the projection
onto their common range, which at this level is \(C_n\). The compatible
finite-dimensional restrictions therefore approximate \(E_{P_{s,t}}\).

\subsection{The angle operator}
Use the inner product \(\langle x,y\rangle=\tr(x^*y)\).
Let \(e_{s,n},e_{t,n}\) be the orthogonal projections of \(L^2(D_n)\)
onto \(B_n(s),B_n(t)\), respectively. Define
\begin{equation}\label{eq:angle}
 T_n=(e_{s,n}e_{t,n}e_{s,n})\big|_{B_n(s)}.
\end{equation}
This is a positive self-adjoint contraction, so its eigenvalues lie
in \([0,1]\).

\begin{lemma}\label{lem:intersection}
The eigenspace of \(T_n\) for eigenvalue \(1\) is exactly \(C_n\).
If \(G_{ij}=\langle f_i,g_j\rangle\), where \((f_i)\) and \((g_j)\)
are orthonormal bases of \(B_n(s)\) and \(B_n(t)\), then the matrix
of \(T_n\) is \(GG^*\).
Moreover, \(S_n=\varnothing\) if and only if \(e_{s,n}\) and
\(e_{t,n}\) commute, where \(S_n\) denotes the eigenvalues in \((0,1)\).
\end{lemma}
\begin{proof}
For \(x\in B_n(s)\),
\(\langle x,T_nx\rangle=\|e_{t,n}x\|_2^2\).
If \(T_nx=x\), equality of the norms of \(x\) and its orthogonal
projection implies \(e_{t,n}x=x\); thus \(x\in C_n\).
Conversely, both projections fix every element of \(C_n\).
Expanding the two projections in the stated orthonormal bases gives
\((GG^*)_{ik}=\sum_j\langle f_i,g_j\rangle\langle g_j,f_k\rangle\).
For the last assertion, put \(e=e_{s,n}\) and \(f=e_{t,n}\).
If \(efe\) has only eigenvalues \(0,1\), it is a projection, and
\[
 0=efe-(efe)^2=ef(1-e)fe=((1-e)fe)^*((1-e)fe).
\]
Hence \((1-e)fe=0\), so \(fe=efe=ef\). Conversely, when \(e,f\)
commute, \(efe=ef\) is a projection.
\end{proof}

It follows that \(c_n\) is the multiplicity of eigenvalue \(1\).
The other data we retain are
\begin{equation}\label{eq:interior}
 S_n=\spec(T_n)\cap(0,1),\qquad
 k_n=|S_n|,\qquad \lambda_n=\max S_n\quad(S_n\ne\varnothing).
\end{equation}
The set \(S_n\) forgets multiplicities. Thus \(k_n\) counts distinct
interior eigenvalues, not eigenvectors. We write a dash for
\(\lambda_n\) when \(S_n=\varnothing\).

\begin{example}\label{ex:lines}
Take two lines spanned by \((1,0)\) and
\((\cos\alpha,\sin\alpha)\) in \(\R^2\), where \(0<\alpha<\pi/2\).
The composition of projections, restricted to the first line, is
multiplication by \(\cos^2\alpha\). This explains the squared cosines
in the tables: an interior eigenvalue \(\lambda\) corresponds to a
principal angle \(\alpha=\arccos\sqrt{\lambda}\). A larger eigenvalue
means a smaller angle. If \(\lambda<1\) is close to \(1\), the two
directions are nearly parallel but still have zero intersection.
\end{example}

The word \emph{interior} in \(S_n\) refers only to the interval \((0,1)\).
These are the spectral angles of the two projections, not the single
interior-angle invariant that also occurs in subfactor theory.
Eigenvalue \(0\) corresponds to a right angle and is deliberately
omitted. In particular, \(S_n=\varnothing\) means that the two
projections commute; it does \emph{not} imply that their ranges coincide.
In the infinite algebra, \emph{commuting} likewise means
\(E_{N_s}E_{N_t}=E_{N_t}E_{N_s}=E_{P_{s,t}}\), not that every element
of \(N_s\) commutes with every element of \(N_t\).

\subsection{A calculation at the first level}\label{sec:first}
\begin{proposition}\label{prop:first}
With multiplicities included, the eigenvalues of \(T_1\) are
\[
 1,\quad1,\quad1,\quad\cos^2(t-s).
\]
Thus \(c_1=4\) if \(t-s\in\pi\Z\), and \(c_1=3\) otherwise.
\end{proposition}
\begin{proof}
In the convention used for the calculation,
\(B_1(r)=u(r)^*\mathcal D_4u(r)\), where \(\mathcal D_4\) is the
algebra of diagonal \(4\times4\) matrices. Its four minimal projections,
multiplied by \(2\), form an orthonormal basis for the normalized trace.
The entries of the cross-Gram matrix are therefore
\(\lvert(u(s)u(t)^*)_{ij}\rvert^2\).
Multiplication of the matrices in \eqref{eq:family}, followed by a
simultaneous reordering of the bases, gives
\[
 G=I_2\oplus
 \begin{pmatrix}
 (1+\cos\Delta)/2&(1-\cos\Delta)/2\\
 (1-\cos\Delta)/2&(1+\cos\Delta)/2
 \end{pmatrix}.
\]
The last block has eigenvalues \(1,\cos\Delta\). Lemma~\ref{lem:intersection}
now gives the assertion.
\end{proof}

For \(\Delta=\pi/2\), the fourth eigenvalue is \(0\). Hence \(c_1=3<4\)
and \(S_1=\varnothing\): this is already an example of a proper
intersection with no interior eigenvalue.

\subsection{Inclusion matrices and finite centers}
The finite-dimensional algebras occurring here are direct sums of
full matrix algebras. Write, with \(M_a(\C)\) denoting all \(a\times a\)
complex matrices,
\[
 C_n\cong\bigoplus_{i=1}^{z_n}M_{a_i}(\C),\qquad
 D_n\cong\bigoplus_{j=1}^{h_n}M_{d_j}(\C).
\]
Then \(c_n=\sum_i a_i^2\) and \(z_n=\dim_{\C}Z(C_n)\).
The \emph{inclusion matrix} \(\Lambda_n\) records how these blocks embed.
Inside the \(j\)-th block of \(D_n\), a suitable basis makes
the action of \(C_n\) block diagonal, with \((\Lambda_n)_{ij}\)
copies of its \(i\)-th matrix block. In particular
\(d_j=\sum_i(\Lambda_n)_{ij}a_i\).
The reported number is
\begin{equation}\label{eq:ell}
 \ell_n=\|\Lambda_n\|^2
       =\text{largest eigenvalue of }\Lambda_n\Lambda_n^{\mathsf T}.
\end{equation}
The norm here is the Euclidean operator norm of a finite matrix.
We compare \(\ell_n\) with \(I(s,t)\), but they are distinct quantities.
For example, below the pair with \(I(s,t)=12\) has \(\ell_1=6\).
No equality at a prescribed finite level is assumed.

Similarly, a finite algebra \(C_n\) may have several central summands
even when its strong limit is a factor. An element central in \(C_n\)
need not commute with elements introduced at later levels.
Center and trace-transition data are recorded in
Appendix~\ref{sec:centers}; they are not used as a proof of factoriality.

\section{The results being illustrated}\label{sec:theorems}
The relevant source is Section~7 of Bakshi--Guin~\cite{BG}.
Their circle parameters are \(a=e^{is}\) and \(b=e^{it}\).
We state the consequences needed here using \(q\) from
\eqref{eq:phase}; equality is separated from the proper-intersection case.
The adjoint convention for the connection conjugates the relative phase,
which does not change \(q\) or the squared cosines.
Let \(e_s,e_t\) be the orthogonal projections onto
\(L^2(N_s),L^2(N_t)\) in \(L^2(M)\). The limiting operator is
\(T=(e_se_te_s)|_{L^2(N_s)}\); its interior spectrum means
\(\spec(T)\cap(0,1)\), with both endpoints excluded.

\Needspace{19\baselineskip}
\begin{theorem}[Bakshi--Guin, in relative-phase notation]\label{thm:bg}
For the above family:
\begin{enumerate}[label=\textup{(\roman*)}]
\item If \(q=1\), then \(N_s=N_t\) and \(I(s,t)=4\).
\item If \(2\leq q<\infty\), then \(P_{s,t}\subset M\) is an
irreducible subfactor and \(I(s,t)=4q\).
\item If \(q=\infty\), then \(I(s,t)=\infty\).
\item For a proper pair, the conditional expectations commute
exactly when \(q=2\).
\item If \(2\leq q<\infty\), the interior spectrum of the limiting
projection product is
\[
 S(q)=\left\{\cos^2\frac{k\pi}{q}:
                   1\leq k\leq\left\lfloor\frac{q-1}{2}\right\rfloor\right\}.
\]
\end{enumerate}
\end{theorem}
The index assertions use \cite[Theorems 7.30, 7.32, 7.36--7.37]{BG};
the commuting and spectral assertions use
\cite[Theorems 7.46--7.47]{BG}. Part (iii) makes no assertion here
about factoriality. \emph{Irreducible} in (ii) means
\(P_{s,t}'\cap M=\C1\): only scalars in \(M\) commute with all of
\(P_{s,t}\). This implies \(Z(P_{s,t})=\C1\).

\paragraph{Why the squared phase appears.}
Adding \(\pi\) to \(t\) swaps rows \(2\) and \(4\) of \(H(t)\).
In this realization this is a relabeling and leaves \(N_t\) unchanged.
Thus \(\theta\) and \(-\theta\) describe the same pair of subfactors.
If \(\Delta/\pi=p/q\) is reduced, the least positive \emph{even}
integer \(m\) with \(\theta^m=1\) is \(2q\). The expression \(2m\)
for the index is therefore \(4q\).
For instance, \(\theta=e^{2\pi i/3}\) has order \(3\), but its
least even exponent giving \(1\) is \(6\); the index is \(12\),
not \(6\). When \(q=1\), repeated row swaps give the equality in (i).

\paragraph{Reading the spectral formula.}
For \(q=2,3,4,5\), respectively, Theorem~\ref{thm:bg} gives
\begin{equation}\label{eq:smallsets}
 \varnothing,\qquad
 \{1/4\},\qquad
 \{1/2\},\qquad
 \left\{\frac{3+\sqrt5}{8},\frac{3-\sqrt5}{8}\right\}.
\end{equation}
For \(q=2\), the empty set is consistent with the presence of a right
angle: its squared cosine \(0\) was excluded from our definition.
The formula gives a \emph{set} and does not give the multiplicity of
each value at a finite level.

The finite subspaces and their orthogonal complements are invariant
under the two projections (they \emph{reduce} them), so
\(S_n\subseteq S_{n+1}\). In the rational cases each \(S_n\) is
contained in \(S(q)\); a value in \(S(q)\) need not appear at the
first level. These observations explain the separate columns
\(k_n\) and \(\lambda_n\).

\begin{remark}
The Fourier-conjugate intersections in Bakshi--Guin--Guruprasad~\cite{BGG}
form a different problem. Their first-level index formula cannot be
substituted for Theorem~\ref{thm:bg}. Proposition~\ref{prop:first}
shows that \(c_1=3\) for every proper pair here, although the indices
of these pairs need not be equal.
\end{remark}

\section{Computations for the Hadamard family}\label{sec:results}
Table~\ref{tab:parameters} lists every tested pair. The labels H01--H14
serve only to locate the same pair in the saved output.
Equality cases were computed through \(n=4\), and all proper pairs
through \(n=5\). Every subsequent table uses the definitions
\[
 c_n=\dim_{\C}C_n,\quad k_n=|S_n|,\quad
 \lambda_n=\max S_n,\quad \ell_n=\|\Lambda_n\|^2.
\]
Recall that \(b_n=4^n\) for this entire section.
Decimal entries are rounded numerical values. A table shared by several
pairs means that all of its displayed entries agree for those pairs.
It does not identify the intersections as embedded algebras.

@@PARAMETERS@@

\subsection[Equality]{Equality: \(t-s\in\pi\Z\)}
We used \((0,0)\), \((\pi/7,\pi/7)\), and
\((\pi/7,8\pi/7)\). The first two compare a matrix with itself,
once at the Fourier member and once elsewhere in the family.
The third compares distinct matrices related by the row swap described
above. In all three cases \(C_n=B_n(s)\), and the computed data are:

@@EQUAL@@

Thus \(c_n=4^n\), every eigenvalue of \(T_n\) is \(1\),
and \(\ell_n=4\). The dash in the \(\lambda_n\) column means that
there is no angle strictly between \(0\) and \(\pi/2\), not that
the calculation is missing.

\subsection[A proper commuting pair]{A proper commuting pair: \(t-s=\pi/2\)}
The pairs \((0,\pi/2)\) and \((\pi/7,9\pi/14)\) have the same
difference and \(q=2\). Theorem~\ref{thm:bg} gives \(I(s,t)=8\)
and commuting conditional expectations.

@@COMMUTING@@

Already \(c_1=3<4=b_1\), so the intersection is proper.
Nevertheless \(S_n=\varnothing\) at every computed level.
This distinguishes the present case from equality without relying
on the interior spectral count alone. The inclusion norm first
equals \(8\) at \(n=2\); its initial value is \(6\).

\subsection[The first noncommuting rational case]{The first noncommuting rational case: \(q=3\)}
We used
\[
 (0,\pi/3),\qquad (0,2\pi/3),\qquad(\pi/7,10\pi/21).
\]
The first and third pairs have difference \(\pi/3\); the second
has difference \(2\pi/3\). Their relative phases have different orders,
but their \emph{squared} phases all have order \(3\).
The predicted index is \(12\), with \(S(3)=\{1/4\}\).

@@ORDER3@@

The value \(1/4\) is present from the first level onward. It corresponds
to the angle \(\arccos(1/2)=\pi/3\), so the conditional expectations
do not commute. The computed \(\ell_n\) first equals \(12\) at
\(n=4\), and has the same value at \(n=5\).
This is agreement with a known index, rather than a deduction of
that index from two equal numerical entries.

\subsection[Larger denominators]{Larger denominators: \(q=4\) and \(q=5\)}
For \((s,t)=(0,\pi/4)\), the predicted index is \(16\)
and \(S(4)=\{1/2\}\).

@@ORDER4@@

The full interior spectral set has appeared, but \(\ell_5=14<16\).
Thus the appearance of all limiting spectral values does not mean
that the inclusion norm has reached the index.

For \((s,t)=(0,\pi/5)\), the predicted index is \(20\).
Its two interior values are approximately \(0.654508\) and
\(0.095492\), as in \eqref{eq:smallsets}.

@@ORDER5@@

The smaller value first appears at \(n=3\). Consequently \(k_n\)
increases from \(1\) to \(2\) while \(\lambda_n\) stays unchanged.
The inclusion norms \(6,8,10,12,14\) are identical to those for \(q=4\);
neither pair has reached its index within the calculated range.

\subsection{Irrational differences}
First consider \((0,1)\) and \((\pi/7,\pi/7+1)\).
Since \(1/\pi\) is irrational, Theorem~\ref{thm:bg} gives infinite
Pimsner--Popa index for both pairs.

@@IRR1@@

The sets through level \(5\), rounded to six decimal places, are
\[
 \begin{aligned}
 S_1=S_2&=\{0.291927\},\\
 S_3=S_4&=\{0.291927,0.173178\},\\
 S_5&=\{0.980085,0.291927,0.173178\}.
 \end{aligned}
\]
The new value at level \(3\) is below the previous maximum, and
the new value at level \(5\) is above it. Looking only at
\(\lambda_1,\ldots,\lambda_4\) would conceal this development.

For the pair \((0,\sqrt2)\), the ratio \(\sqrt2/\pi\) is also
irrational: a nonzero rational value for this ratio would imply
that \(\pi\) is algebraic. The index is again infinite.

@@IRR2@@

Here \(S_5\), rounded, is
\(\{0.905092,0.204903,0.024318\}\).
The dimension and inclusion-norm sequences match the previous
irrational pair, but the spectral locations differ. The first
increase in the largest interior eigenvalue occurs at level \(3\)
rather than at level \(5\).

\subsection{A small perturbation of a rational difference}
Finally, take
\[
 (s,t)=(0,\pi/3+1/100).
\]
This changes the index-\(12\) example by \(0.01\) radians.
Its ratio \(\Delta/\pi=1/3+1/(100\pi)\) is irrational, since
rationality would force \(\pi\) to be rational. Its theoretical
index is therefore infinite.

@@PERTURB@@

At level \(5\), this pair has \(c_5=462\), whereas the unperturbed
pair has \(c_5=463\). The perturbed pair also has an interior
eigenvalue approximately \(0.999100\). It must not be counted as \(1\).
An eigenvector belonging to such an eigenvalue is orthogonal to the
eigenvalue-\(1\) space \(C_5\); it describes nearly parallel
directions outside the intersection, as in Example~\ref{ex:lines}.
Thus closeness of the eigenvalue to \(1\) does not mean closeness
of its eigenvector to the intersection.

\subsection{What these comparisons establish}
The three pairs of common shifts
\[
 \begin{gathered}
 (0,\pi/2),\ (\pi/7,9\pi/14);\\
 (0,\pi/3),\ (\pi/7,10\pi/21);\\
 (0,1),\ (\pi/7,\pi/7+1)
 \end{gathered}
\]
agree in the calculated dimensions, simple block sizes, center
dimensions, inclusion norms, and interior spectra with multiplicities.
The same comparisons agree for \((0,\pi/3)\) and \((0,2\pi/3)\).
These are numerical comparisons over the stated levels.

The \(q=4\), \(q=5\), and all irrational examples share
\[
 (c_1,\ldots,c_5)=(3,10,35,126,462),\qquad
 (\ell_1,\ldots,\ell_5)=(6,8,10,12,14).
\]
Nevertheless their theoretical indices are \(16\), \(20\), and
\(\infty\), respectively. This is the principal limitation of
reading off a limiting index from a short sequence.
Likewise, finite-level centers do not determine whether the limit
is a factor. The rational/irrational classification comes from
Theorem~\ref{thm:bg}, not from extrapolating these lists.

\appendix
\section[The separate S(1,1,2,2) family]{The separate \(S(1,1,2,2)\) family}\label{sec:ah}
The demo also retains a connection family associated with the four-arm
graph \(S(1,1,2,2)\), whose arm lengths are \(1,1,2,2\); see
Bisch--C\'aceres~\cite{BC} for the construction.
Denote its subfactors by \(\widetilde N_r\subset\widetilde M\)
to distinguish them from the Hadamard family. Their individual index is
\[
 [\widetilde M:\widetilde N_r]=\frac{5+\sqrt{17}}2
 \simeq4.561553.
\]
We use the definitions \(b_n,c_n,k_n,\lambda_n,\ell_n\) from the
main text for this separate tower. The Hadamard theorem does not
give the index of \(\widetilde N_s\cap\widetilde N_t\).
Throughout this appendix \(s=\pi/3\).

For \(t=s\), the intersection equals \(\widetilde N_s\).
The computation gives:

@@AH_EQUAL@@

For all seven unequal parameters listed in Table~\ref{tab:ahmax},
the following displayed data agree.

@@AH_PROPER@@

@@AH_MAXIMA@@

Every pair in the latter table has a specific interpretation.
The pair with \(t=\pi/3\) is equality. For \(t=\pi/2\) and \(t=7/5\),
the largest interior eigenvalue remains constant through all four
levels, even though the number of distinct interior eigenvalues grows.
The value \(t=2\pi/3=\pi-s\) also gives a constant maximum, but does
not give the same embedded subfactor as \(s\).
For \(t=\pi\) and \(t=4\pi/3\), the maximum increases at level \(3\).
In particular, \(t-s=\pi\) in the latter example does not imply
equality in this family. For \(t=5\pi/3=2\pi-s\) the maximum
increases at each displayed step. For \(t=1/10\) it increases
slightly at level \(3\).

The computed inclusions of the intersection algebras have two connected
components through these four levels. A connected component groups simple
summands linked by nonzero inclusion multiplicities. Its sum of central
projections has trace approximately \(0.179806\) for one component and
\(0.820194\) for the other. The initial vertex labels in the program
are \(d2\) and \(\{A,a2\}\), respectively; the labels have no further
role in interpreting these trace values.
If these compatible projections remained central at every subsequent
level, they would establish that the limiting intersection is not a
factor. Four levels alone do not prove that persistence. Nor does the
observed growth of \(k_n\) or \(\ell_n\) prove infinite index.

\section{Center and trace-transition data}\label{sec:centers}
For completeness, we record the remaining columns of the computation
and define the transition statistic explicitly. Let \(L^{(n)}\) be
the inclusion matrix for \(C_n\subset C_{n+1}\); it is a different
matrix from \(\Lambda_n\), which describes \(C_n\subset D_n\).
Let \(w_i^{(n)}\) be the trace of a minimal projection in the
\(i\)-th matrix summand of \(C_n\). Define
\[
 K^{(n)}_{ij}
 =\frac{L^{(n)}_{ij}w_j^{(n+1)}}{w_i^{(n)}},
 \qquad K_{1,n}=K^{(1)}K^{(2)}\cdots K^{(n-1)}.
\]
Each row of \(K^{(n)}\) sums to \(1\), by compatibility of the
trace with the inclusion. Thus its rows can be viewed as probability
distributions on the next level's summands. The recorded coefficient is
\[
 d_{1,n}=\frac12\max_{i,i'}
              \sum_j |(K_{1,n})_{ij}-(K_{1,n})_{i'j}|,
 \qquad n\geq2.
\]
It lies in \([0,1]\). A value \(0\) means all these distributions
agree; a value \(1\) means that two have disjoint supports.
There is no transition at \(n=1\), so the output puts a dash there.

@@CENTERS@@

For the equal Hadamard pairs, \(z_n\) alternates between \(4\) and
\(1\), although the limit is the known factor \(N_s\).
For the unequal \(S(1,1,2,2)\) pairs, \(z_4=7\), although the
four-level diagram has two connected components: one component
can contain several simple summands.

The coefficient \(d_{1,n}\) compares only distributions starting at
the first level. It does not, by itself, control information introduced
at later levels. Its vanishing, or its decrease over a finite range,
is therefore not used here as a factoriality criterion.

\section{Numerical procedure and reproducibility}\label{sec:numerics}
The calculations use double-precision complex arithmetic and
orthonormal bases for the trace inner product. Lemma~\ref{lem:intersection}
reduces the spectral calculation to the squared singular values
of a cross-Gram matrix. The implementation uses tolerance
\(\varepsilon=10^{-9}\): an eigenvalue at least
\((1-\varepsilon)^2\) is assigned to the intersection.
Interior eigenvalues are grouped with tolerance \(10^{-8}\);
the comparisons use \(5\cdot10^{-7}\).
These procedures do not constitute rigorous error bounds.
For the perturbed example, the observed gap \(1-0.999100\)
is much larger than the intersection cutoff.

The saved run contains \(22\) pairs and @@ROWS@@ finite-level records.
All @@CHECKS@@ per-pair checks and four cross-pair comparisons passed.
The checks include Proposition~\ref{prop:first}, the equality cases,
containment in the rational spectral sets, nesting of the interior
spectra, and \(c_n=\sum_i a_i^2\).
For rational phases, the computed inclusion norms were also checked
not to exceed the known index. The three existing tests for equality,
the Fourier-conjugate example, and nested spectra passed.
The tables in this note are generated from that saved run; no new
numerical run was needed for this exposition.

The parameters are specified exactly before evaluating matrices
numerically. In particular, the distinction between a rational
and an irrational multiple of \(\pi\) is not inferred by rounding
a decimal approximation.

The source files are \path{Code/intersection_demo.py},
\path{Code/intersection_demo_results.json}, and
\path{Code/intersection_demo_output.txt}.
The prose is in \path{Code/intersection_report_template.tex};
\path{Code/write_intersection_report.py} inserts the recorded tables
and emits a standalone \LaTeX{} source. From the project root:
\begin{quote}\small\ttfamily
.\textbackslash.venv\textbackslash Scripts\textbackslash python.exe Code/intersection\_demo.py\\
.\textbackslash.venv\textbackslash Scripts\textbackslash python.exe Code/write\_intersection\_report.py\\
pdflatex -output-directory=output/pdf\\
\hspace*{1em}output/pdf/intersection\_results\_explained.tex
\end{quote}
The last two printed lines form one command. Run that command twice
to resolve references. The generator expects the complete default
calculation. The standalone generated source includes every table
and does not need the data file when compiled.

\begin{thebibliography}{9}
\bibitem{BG}
K.~C. Bakshi and S. Guin,
\emph{Relative position between a pair of spin model subfactors},
\href{https://arxiv.org/html/2212.07197v3}{arXiv:2212.07197v3},
Section~7. We specify this version for its theorem numbering;
the second part was published separately as
\emph{A family of subfactors arising from a pair of complex Hadamard matrices}.
\bibitem{BGG}
K.~C. Bakshi, S. Guin and Guruprasad,
\emph{Intersection of conjugate Hadamard subfactors arising from Fourier matrices},
\href{https://arxiv.org/abs/2510.14499}{arXiv:2510.14499}.
\bibitem{BC}
D. Bisch and J. C\'aceres,
\emph{New hyperfinite subfactors with infinite depth},
\href{https://arxiv.org/abs/2509.00227}{arXiv:2509.00227}.
\end{thebibliography}
\end{document}
